% GENERATED FILE. Edit canonical lessons, then run npm run generate:books.
% canonical-source-hash: fnv1a64:29b1ef0bf7ade706
% canonical-lessons: MR-C416-khunti, MR-C416-dhobi, MR-C416-pilne, MR-C416-hamal, MR-C416-simpadne

\chapter{Washermen and Porters}
\label{ch:mr-washermen-and-porters}
\hlchaptermodality{416}

\begin{chapteropening}
\textbf{By the end of this chapter:} \emph{I can talk about pegs, washermen and porters, and wringing and sprinkling clothes.}

Complete the last lesson of chapter 416: I can talk about pegs, washermen and porters, and wringing and sprinkling clothes.
\end{chapteropening}

\section[\texorpdfstring{\mr{खुंटी} (khuṇṭī)}{खुंटी (khuṇṭī)}]{\mr{खुंटी} (khuṇṭī) — a peg; a hook}

\label{lesson:MR-C416-khunti}



\begin{quote}
Before the new one: say the Marathi for to spread out bedding, then the Marathi for to cover oneself.
\end{quote}

\subsection*{You'll want to know: \mr{खुंटी}}
\textbf{\mr{खुंटी}} — \emph{khuṇṭī} — \textquotedblleft{}a peg; a hook\textquotedblright{}.

\begin{quote}
\mr{दाराच्या} \mr{मागे} \mr{खुंटीवर} \mr{पिशवी} \mr{टांगली} \mr{आहे}.
\end{quote}

\begin{quote}
\emph{dārācyā māge khuṇṭīvar piśvī ṭāṅglī āhe.}
\end{quote}

\begin{quote}
A bag is hanging on a peg behind the door.
\end{quote}

Feminine.

A peg in the wall for clothes and bags; \mr{खुंटीवर}, on the peg.

\begin{cousinweb}
\textbf{\mr{खुंटी}} (\emph{khuṇṭī}) is a Marathi word for a wall peg; its origin is not settled.
\end{cousinweb}

\begin{grammarlens}[title={What you've built}]
One more word for getting around.
\end{grammarlens}

\subsection*{Your turn}
\begin{itemize}
\raggedright
  \item \emph{Say it:} \emph{khuṇṭī}
  \item \emph{Say it:} \emph{khuṇṭī}, once more
  \item \emph{Say it:} \emph{dārācyā māge khuṇṭīvar piśvī ṭāṅglī āhe.} — A bag is hanging on a peg behind the door.
  \item \emph{Recall:} say the Marathi for to spread out bedding, then the Marathi for to cover oneself, then say \emph{khuṇṭī} again
\end{itemize}

\begin{tcolorbox}[breakable,colback=teal!4,colframe=teal!35!black,title={Before you move on}]
What does \mr{खुंटी} mean? (\textquotedblleft{}A peg; a hook\textquotedblright{}.) Say it once more.
\end{tcolorbox}
\section[\texorpdfstring{\mr{धोबी} (dhobī)}{धोबी (dhobī)}]{\mr{धोबी} (dhobī) — a washerman}

\label{lesson:MR-C416-dhobi}



\begin{quote}
Before the new one: say the Marathi for to cover oneself, then the Marathi for a peg.
\end{quote}

\subsection*{You'll want to know: \mr{धोबी}}
\textbf{\mr{धोबी}} — \emph{dhobī} — \textquotedblleft{}a washerman\textquotedblright{}.

\begin{quote}
\mr{धोबी} \mr{दर} \mr{आठवड्याला} \mr{कपडे} \mr{नेतो}.
\end{quote}

\begin{quote}
\emph{dhobī dar āṭhvaḍyālā kapḍe neto.}
\end{quote}

\begin{quote}
The washerman takes the clothes every week.
\end{quote}

Masculine.

The Marathi word \mr{परीट} names the same trade.

\begin{cousinweb}
\textbf{\mr{धोबी}} (\emph{dhobī}) is a loan from Hindi \emph{dhobī}, \textquotedblleft{}washerman\textquotedblright{}: CDIAL lists it under \emph{*dhauvati}, \textquotedblleft{}washes\textquotedblright{}, and notes that the Hindi word spread to the neighbouring languages, Marathi among them (CDIAL 6886).
\end{cousinweb}

\begin{grammarlens}[title={What you've built}]
One more word for getting around.
\end{grammarlens}

\subsection*{Your turn}
\begin{itemize}
\raggedright
  \item \emph{Say it:} \emph{dhobī}
  \item \emph{Say it:} \emph{dhobī}, once more
  \item \emph{Say it:} \emph{dhobī dar āṭhvaḍyālā kapḍe neto.} — The washerman takes the clothes every week.
  \item \emph{Recall:} say the Marathi for to cover oneself, then the Marathi for a peg, then say \emph{dhobī} again
\end{itemize}

\begin{tcolorbox}[breakable,colback=teal!4,colframe=teal!35!black,title={Before you move on}]
What does \mr{धोबी} mean? (\textquotedblleft{}A washerman\textquotedblright{}.) Say it once more.
\end{tcolorbox}
\section[\texorpdfstring{\mr{पिळणे} (piḷṇe)}{पिळणे (piḷṇe)}]{\mr{पिळणे} (piḷṇe) — to wring}

\label{lesson:MR-C416-pilne}



\begin{quote}
Before the new one: say the Marathi for a peg, then the Marathi for a washerman.
\end{quote}

\subsection*{You'll want to know: \mr{पिळणे}}
\textbf{\mr{पिळणे}} — \emph{piḷṇe} — \textquotedblleft{}to wring\textquotedblright{}.

\begin{quote}
\mr{तिने} \mr{ओला} \mr{सदरा} \mr{पिळला}.
\end{quote}

\begin{quote}
\emph{tine olā sadrā piḷlā.}
\end{quote}

\begin{quote}
She wrung out the wet shirt.
\end{quote}

Past, with -\mr{ने} (he …): \mr{त्याने} \mr{पिळलं}, \emph{tyāne piḷlaṁ}.

Also for fruit: \mr{लिंबू} \mr{पिळणे}, to squeeze a lemon.

\begin{cousinweb}
\textbf{\mr{पिळणे}} (\emph{piḷṇe}) is inherited: CDIAL lists Marathi \emph{piḷṇẽ}, \textquotedblleft{}to squeeze\textquotedblright{}, under Sanskrit \emph{pīḍayati}, \textquotedblleft{}presses, squeezes\textquotedblright{} (CDIAL 8226).
\end{cousinweb}

\begin{grammarlens}[title={What you've built}]
One more word for getting around.
\end{grammarlens}

\subsection*{Your turn}
\begin{itemize}
\raggedright
  \item \emph{Say it:} \emph{piḷṇe}
  \item \emph{Say it:} \emph{piḷṇe}, once more
  \item \emph{Say it:} \emph{tine olā sadrā piḷlā.} — She wrung out the wet shirt.
  \item \emph{Recall:} say the Marathi for a peg, then the Marathi for a washerman, then say \emph{piḷṇe} again
\end{itemize}

\begin{tcolorbox}[breakable,colback=teal!4,colframe=teal!35!black,title={Before you move on}]
What does \mr{पिळणे} mean? (\textquotedblleft{}To wring\textquotedblright{}.) Say it once more.
\end{tcolorbox}
\section[\texorpdfstring{\mr{हमाल} (hamāl)}{हमाल (hamāl)}]{\mr{हमाल} (hamāl) — a porter}

\label{lesson:MR-C416-hamal}



\begin{quote}
Before the new one: say the Marathi for a washerman, then the Marathi for to wring.
\end{quote}

\subsection*{You'll want to know: \mr{हमाल}}
\textbf{\mr{हमाल}} — \emph{hamāl} — \textquotedblleft{}a porter\textquotedblright{}.

\begin{quote}
\mr{स्टेशनवर} \mr{हमालाने} \mr{आमचं} \mr{सामान} \mr{उचललं}.
\end{quote}

\begin{quote}
\emph{sṭeśanvar hamālāne āmcaṁ sāmān ucallaṁ.}
\end{quote}

\begin{quote}
At the station a porter lifted our luggage.
\end{quote}

Masculine.

At stations they wear a red shirt and a brass badge.

\begin{cousinweb}
\textbf{\mr{हमाल}} (\emph{hamāl}) is a loan through Persian, usually traced to Arabic \emph{ḥammāl}, \textquotedblleft{}a carrier of loads\textquotedblright{}.
\end{cousinweb}

\begin{grammarlens}[title={What you've built}]
One more word for getting around.
\end{grammarlens}

\subsection*{Your turn}
\begin{itemize}
\raggedright
  \item \emph{Say it:} \emph{hamāl}
  \item \emph{Say it:} \emph{hamāl}, once more
  \item \emph{Say it:} \emph{sṭeśanvar hamālāne āmcaṁ sāmān ucallaṁ.} — At the station a porter lifted our luggage.
  \item \emph{Recall:} say the Marathi for a washerman, then the Marathi for to wring, then say \emph{hamāl} again
\end{itemize}

\begin{tcolorbox}[breakable,colback=teal!4,colframe=teal!35!black,title={Before you move on}]
What does \mr{हमाल} mean? (\textquotedblleft{}A porter\textquotedblright{}.) Say it once more.
\end{tcolorbox}
\section[\texorpdfstring{\mr{शिंपडणे} (śimpaḍṇe)}{शिंपडणे (śimpaḍṇe)}]{\mr{शिंपडणे} (śimpaḍṇe) — to sprinkle}

\label{lesson:MR-C416-simpadne}



\begin{quote}
Before the new one: say the Marathi for to wring, then the Marathi for a porter.
\end{quote}

\subsection*{You'll want to know: \mr{शिंपडणे}}
\textbf{\mr{शिंपडणे}} — \emph{śimpaḍṇe} — \textquotedblleft{}to sprinkle\textquotedblright{}.

\begin{quote}
\mr{इस्त्रीआधी} \mr{तो} \mr{कपड्यांवर} \mr{पाणी} \mr{शिंपडतो}.
\end{quote}

\begin{quote}
\emph{istrīādhī to kapaḍyāṁvar pāṇī śimpaḍto.}
\end{quote}

\begin{quote}
Before ironing he sprinkles water on the clothes.
\end{quote}

Past, with -\mr{ने} (he …): \mr{त्याने} \mr{शिंपडलं}, \emph{tyāne śimpaḍlaṁ}.

The place sprinkled takes -\mr{वर}: \mr{अंगणात} \mr{पाणी} \mr{शिंपडणे}, to sprinkle water in the yard.

\begin{cousinweb}
\textbf{\mr{शिंपडणे}} (\emph{śimpaḍṇe}) is inherited: CDIAL lists Marathi \emph{śĩpaḍṇẽ}, \textquotedblleft{}to sprinkle, scatter\textquotedblright{}, under Sanskrit \emph{sikta}, \textquotedblleft{}sprinkled, wetted\textquotedblright{} (CDIAL 13388).
\end{cousinweb}

\begin{grammarlens}[title={What you've built}]
One more word for getting around.
\end{grammarlens}

\subsection*{Your turn}
\begin{itemize}
\raggedright
  \item \emph{Say it:} \emph{śimpaḍṇe}
  \item \emph{Say it:} \emph{śimpaḍṇe}, once more
  \item \emph{Say it:} \emph{istrīādhī to kapaḍyāṁvar pāṇī śimpaḍto.} — Before ironing he sprinkles water on the clothes.
  \item \emph{Recall:} say the Marathi for to wring, then the Marathi for a porter, then say \emph{śimpaḍṇe} again
\end{itemize}

\begin{tcolorbox}[breakable,colback=teal!4,colframe=teal!35!black,title={Before you move on}]
What does \mr{शिंपडणे} mean? (\textquotedblleft{}To sprinkle\textquotedblright{}.) Say it once more.
\end{tcolorbox}
